% ------------------------------------------------------------------------------
% SETUP ------------------------------------------------------------------------
% ------------------------------------------------------------------------------

\documentclass[12pt]{article}
\usepackage[a4paper, width = 175mm, top = 20mm, bottom = 20mm, 
bindingoffset = 0mm]{geometry}
\usepackage[utf8]{inputenc}
\usepackage{url}
\usepackage{ragged2e}
\usepackage{array}
\usepackage{tabularx}
\usepackage{float}
\usepackage{xcolor, listings}
\usepackage{textcomp}
\usepackage{fancyhdr}
\newcommand{\changefont}{%
    \fontsize{8}{11}\selectfont
}

\usepackage[hang, flushmargin, bottom]{footmisc} 
\usepackage{hyperref}
\hypersetup{
  colorlinks = true,
  linkcolor = .,
  urlcolor = .,
  citecolor = .,
  bookmarks = true}

\pagestyle{fancy}
\fancyhead{}
\fancyhead[R]{\changefont{Topic-specific sentiment analysis for tweets by 
German MPs}}
\fancyfoot{}
\fancyfoot[R]{\thepage}
\setlength{\headheight}{14.5pt}
\setlength{\parindent}{0pt}
\interfootnotelinepenalty = 10000

% ------------------------------------------------------------------------------
% MAIN -------------------------------------------------------------------------
% ------------------------------------------------------------------------------

\usepackage{/usr/share/R/share/texmf/tex/latex/Sweave}
\begin{document}
\Sconcordance{concordance:einladung.tex:einladung.Rnw:%
1 43 1 1 0 94 1}


% FRONT PAGE -------------------------------------------------------------------
 
\begin{titlepage}
\begin{center}
    
\large
Statistical Consulting
    
\vspace{0.5cm}
      
\rule{\textwidth}{1pt}
% \large
\textbf{Topic-specific sentiment analysis for tweets by German MPs}
\rule{\textwidth}{1pt}
   
\vspace{0.5cm}
      
\normalsize
Department of Statistics \\
Ludwig-Maximilians-Universität München 

\vspace{0.5cm}

\includegraphics[width = 0.2\textwidth]{figures/sigillum.png}

\vspace{0.5cm}

\rule{\textwidth}{1pt}
\small
% \raggedright
New data sources offer vast potential for studying public opinion in the 
political environment.
Frequently arising research questions concern the nature of topics discussed 
and the sentiment expressed toward them.
We perform sentiment analysis for tweets by German MPs issued since the last 
federal election in September 2017, scraped from the Web, with special regard to 
topical context.
The contribution of this work is twofold.
First, we explore how such analysis can be conducted with either standard 
machine learning approaches for tabular data or more complex, BERT-based deep 
learning solutions.
Our results suggest that casting the problem at hand as a standard machine 
learning task requires extensive feature engineering and achieves only moderate
predictive performance, whereas BERT-based models yield satisfactory 
results for sentiment analysis.
We further note that topic modeling does not aid classification in this 
application but attribute this largely to the difficulty of topic extraction.
Second, we provide teaching material for instructing fellow researchers in the 
applied techniques.
With this collection, which will remain publicly available, we hope to 
facilitate the entrance to NLP for practitioners from other disciplines.
\rule{\textwidth}{1pt}

% \center

\vspace{0.3cm}

\begin{tabular}{rl}
  Speakers & Asmik Nalmpatian \& Lisa Wimmer \\
  & \\
  Project partner & Prof. Dr. Paul Thurner \\
  & \textit{Department of Political Science} \\
  & \\
  Supervisors & Matthias Aßenmacher, Ph.D. \\
  & Prof. Dr. Christian Heumann \\
  & \textit{Department of Statistics} \\
  & \\
  Date & Monday, July 12\textsuperscript{th}, 2021, 16:00 s.t. \\
  & \\
  Meeting & ID 918 5472 5495, passcode 413775
\end{tabular}

\vspace{0.5cm}

\rule{\textwidth}{1pt}

\vspace{0.5cm}

\url{https://lmu-munich.zoom.us/j/91854725495?pwd=NHVLemt4ZzZHNk8vSUpHWHBZamZjUT09}

\end{center}
\end{titlepage}

\end{document}
